\section{Pengumpulan Data}

\subsection{Sumber Data}

Sumber data utama dalam penelitian ini mengacu pada skema basis data \texttt{rentals} yang telah dirancang pada sistem. Mengingat sistem ini berada pada tahap pra-peluncuran (\textit{pre-launching}) dan belum memiliki data historis riil dalam jumlah yang cukup untuk pelatihan \textit{machine learning}, penelitian ini menggunakan dataset sintetis. Karakteristik, distribusi, dan anomali data disesuaikan dengan {\bf kondisi faktual} di Rusunami The Jarrdin Cihampelas (RTJC) berdasarkan hasil wawancara dengan agen dan pengelola.

Dataset ini terdiri dari atribut-atribut yang merepresentasikan data penyewa dan perilaku transaksi, sebagaimana didefinisikan dalam struktur tabel \texttt{rentals} berikut.

\begin{itemize}
    \item \textbf{Identitas penyewa:} atribut \texttt{nama} dan \texttt{nik} (Nomor Induk Kependudukan) untuk identifikasi unik.
    \item \textbf{Profil demografis:}
    \begin{itemize}
        \item \texttt{status\_pasutri}: menyatakan status pernikahan ('Menikah' atau 'Belum Menikah'), yang krusial untuk mendeteksi potensi pelanggaran asusila pada penyewaan satu unit.
        \item \texttt{status\_kewarganegaraan}: Membedakan antara 'WNI' dan 'WNA' untuk keperluan pengawasan orang asing.
    \end{itemize}
    \item \textbf{Detail transaksi sewa:}
    \begin{itemize}
        \item \texttt{jenis\_sewa}: kategori durasi sewa ('Harian', 'Mingguan', 'Bulanan').
        \item \texttt{unit\_number}: unit apartemen yang disewa.
        \item \texttt{metode\_pembayaran}: metode pembayaran yang dipilih penyewa, seperti 'Cash', 'Transfer', 'Kartu Kredit', dll. Metode pembayaran tunai (\textit{Cash}) sering kali menjadi indikator penyewa berpotensi melanggar karena transaksi yang lebih sulit dilacak.
    \end{itemize}
    \item \textbf{Data temporal (waktu):}
    \begin{itemize}
        \item \texttt{tanggal\_checkin} dan \texttt{waktu\_checkin}: penanda waktu masuk.
        \item \texttt{tanggal\_checkout} dan \texttt{waktu\_checkout}: penanda waktu keluar.
        \item \texttt{lama\_menginap}: durasi lama inap yang dihitung secara otomatis.
    \end{itemize}
    \item \textbf{Atribut target (labeling):}
    \begin{itemize}
        \item \texttt{komentar}: kolom teks yang berisi catatan pelanggaran(jika ada). 
        \item \texttt{classification}: atribut klasifikasi bertipe \textit{enum} ('normal', 'tidak normal') yang berfungsi sebagai label kelas target untuk pelatihan model, diturunkan dari isi kolom komentar.
    \end{itemize}
\end{itemize}

Proses pembuatan data dilakukan menggunakan pustaka \texttt{Faker} untuk menjamin variasi data identitas, namun tetap terikat pada keadaan atau kasus yang ditemukan di lapangan (contoh: penyewa 'transit' cenderung memiliki durasi sewa sangat pendek dan melakukan \textit{check-in} di jam-jam tertentu).
